% ==================== 1.4 论文主要研究内容及章节安排 ====================
\section{论文主要研究内容及章节安排}
\label{sec:content_arrangement}

本节介绍论文的主要研究内容、技术路线以及各章节安排。

\subsection{研究内容与关键技术路线}
\label{subsec:research_content}

针对三大科学问题，本文开展以下四个方面的研究工作：

\textbf{（1）月面巡视器自主行走数字孪生框架构建}

在分析经典数字孪生五维模型适用性的基础上，针对月面巡视器自主行走场景的特殊需求，提出增强的五维数字孪生模型，从物理实体、虚拟实体、服务、孪生数据、连接五个维度进行适应性设计。在此基础上，构建月面巡视器自主行走数字孪生的总体架构，明确各模块功能与接口关系，为后续研究提供统一的框架支撑。

\textbf{（2）月面环境与巡视器动力学数字孪生建模}

针对月面环境建模需求，研究月面地形、月壤环境、光照温度场的数字孪生建模方法。基于DEM数据和高程信息，建立月面地形数字孪生模型；基于离散元方法，建立月壤力学特性模型；考虑太阳运动规律，建立月面光照与温度场模型。

针对巡视器动力学建模需求，研究轮-壤交互力学理论和多体动力学建模方法。建立考虑滑移特性的轮-壤相互作用模型，构建巡视器多体动力学数字孪生模型，并通过仿真验证模型的准确性。

\textbf{（3）基于数字孪生的月面岩石识别与避障}

研究数字孪生平台下的月面岩石识别方法。构建岩石识别框架，开发增强型ORB-SLAM2系统，实现对月面岩石目标的精确识别与定位。在此基础上，设计基于岩石识别结果的避障策略，提高巡视器在复杂地形环境下的行走安全性。

\textbf{（4）基于数字孪生的巡视器路径规划}

研究数字孪生环境下的路径规划方法。构建数字孪生路径规划框架，提出基于A-D3QN（Attention-based Double Deep Q-Network）的混合路径规划算法，实现全局规划与局部规划的协同优化。通过仿真实验验证算法的有效性和优越性。

本文的技术路线如图\ref{fig:roadmap}所示。

% 技术路线图（示意）
\begin{figure}[htbp]
    \centering
    \begin{tikzpicture}[
        node distance=1.5cm,
        block/.style={rectangle, draw, fill=blue!10, text width=6cm, text centered, minimum height=1cm, rounded corners},
        arrow/.style={->, >=stealth, thick}
    ]
        \node[block] (p1) {数字孪生框架构建\\增强五维模型、总体架构};
        \node[block, below=of p1] (p2) {环境数字孪生建模\\地形、月壤、光照温度};
        \node[block, below=of p2] (p3) {动力学数字孪生建模\\轮-壤交互、多体动力学};
        \node[block, below=of p3] (p4) {岩石识别与避障\\增强ORB-SLAM2、避障策略};
        \node[block, below=of p4] (p5) {路径规划\\A-D3QN混合规划算法};
        
        \draw[arrow] (p1) -- (p2);
        \draw[arrow] (p2) -- (p3);
        \draw[arrow] (p3) -- (p4);
        \draw[arrow] (p4) -- (p5);
    \end{tikzpicture}
    \caption{论文技术路线图}
    \label{fig:roadmap}
\end{figure}

\subsection{本文章节安排}
\label{subsec:chapter_arrangement}

本文共分七章，各章内容安排如下：

\textbf{第1章 绪论}：阐述研究背景和意义，综述国内外研究现状，分析现存难点和问题，提出研究内容和章节安排。

\textbf{第2章 基于增强五维模型的月面巡视器自主行走数字孪生框架}：分析经典五维模型的局限性，提出面向月面巡视器自主行走的增强五维模型，构建数字孪生总体架构。

\textbf{第3章 月面环境数字孪生建模与仿真}：研究月面地形、月壤环境、光照温度场的数字孪生建模方法，进行实验验证和结果分析。

\textbf{第4章 月面巡视器动力学数字孪生建模}：研究轮-壤交互力学理论和多体动力学建模方法，构建巡视器动力学数字孪生模型。

\textbf{第5章 基于数字孪生的月面岩石识别与避障}：构建岩石识别框架，开发增强型ORB-SLAM2系统，设计避障策略。

\textbf{第6章 基于数字孪生的巡视器路径规划}：构建路径规划框架，提出A-D3QN混合路径规划算法，进行仿真验证。

\textbf{第7章 总结与展望}：总结全文研究工作，凝练主要创新点，展望未来研究方向。
